\par
Concluimos considerando la «representabilidad» de un grupoide
cofibrado. Sea $\underline C$ una categoría fija. Todo objeto $R$ de
$\underline C$ define un funtor
\[
 h_R:\underline C\longrightarrow(\mathrm{Sets}),
 \qquad h_R(S)=\operatorname{Hom}_{\underline C}(R,S),
\]
y, por tanto, un funtor plenamente fiel
\[
 h:\underline C^{\mathrm{op}}
 \longrightarrow
 \underline{\operatorname{Hom}}(\underline C,(\mathrm{Sets})).
\]
Una \emph{cocategoría en $\underline C$} es un diagrama
\[
\begin{tikzcd}[column sep=huge]
 h_{R_1}\times_{h_{R_0}}h_{R_1} \arrow[r]
 & h_{R_1}
   \arrow[r,shift left=3,"{s_1}"]
   \arrow[r,shift right=3,"{s_2}"']
 & h_{R_0} \arrow[l,"{\epsilon}" description]
\end{tikzcd}
\qquad
\text{con }s_i\circ\epsilon=\operatorname{id}\quad(i=1,2),
\]
donde la primera flecha se denota por $c$ y
\[
 h_{R_1}\times_{h_{R_0}}h_{R_1}
 = (h_{R_1})_{s_1}\times_{h_{R_0}}{}_{s_2}(h_{R_1}),
\]
tal que

(i) $c$ es asociativo y compatible con $s_i$ para $i=1,2$;

(ii) ambos compuestos
\begin{align*}
 h_{R_1}&\xrightarrow{\ \sim\ }
 h_{R_0}\times_{h_{R_0}}h_{R_1}
 \xrightarrow{\ \epsilon\times1\ }
 h_{R_1}\times_{h_{R_0}}h_{R_1}
 \xrightarrow{\ c\ }h_{R_1},\\[1ex]
 h_{R_1}&\xrightarrow{\ \sim\ }
 h_{R_1}\times_{h_{R_0}}h_{R_0}
 \xrightarrow{\ 1\times\epsilon\ }
 h_{R_1}\times_{h_{R_0}}h_{R_1}
 \xrightarrow{\ c\ }h_{R_1}
\end{align*}
son identidades.

Si $(R_0,R_1,s_1,s_2,\epsilon,c)$ es una cocategoría en
$\underline C$, para todo $S\in\underline C$ definimos una categoría
$h_{(R_0,R_1,\ldots)}(S)$ mediante
\[
 \operatorname{ob}h_{(R_0,R_1,\ldots)}(S)
   =\operatorname{Hom}_{\underline C}(R_0,S),
 \qquad
 \operatorname{Fl}h_{(R_0,R_1,\ldots)}(S)
   =\operatorname{Hom}_{\underline C}(R_1,S),
\]
con la composición dada por
\[
 c:h_{R_1}\times_{h_{R_0}}h_{R_1}\longrightarrow h_{R_1}.
\]
Las categorías así obtenidas dependen funtorialmente de $S$ y definen
una categoría cofibrada $\underline h_{(R_0,R_1,\ldots)}$ sobre
$\underline C$. Por ejemplo, todo objeto $R$ de $\underline C$ define
una cocategoría discreta
\[
 \underline h_R
 =\underline h_{(R,R,s_1,s_2,\epsilon,c)}
\]
tomando $s_1,s_2,\epsilon$ y $c$ todos iguales a los morfismos
identidad. Un homomorfismo
\[
 (R_0,R_1,\ldots)\longrightarrow(R'_0,R'_1,\ldots)
\]

%%%%%%%%%%  scan idx 72  |  folio 61  %%%%%%%%%%

de cocategorías es un par de morfismos de funtores representados
$h_{R_0}\longrightarrow h_{R'_0}$ y
$h_{R_1}\longrightarrow h_{R'_1}$ cuyo morfismo inducido
\[
 h_{(R_0,R_1,\ldots)}(S)
 \longrightarrow h_{(R'_0,R'_1,\ldots)}(S)
\]
es un funtor para todo $S\in\underline C$. Toda cocategoría
$(R_0,R_1,\ldots)$ posee, en particular, un homomorfismo
\[
 \epsilon:\underline h_{R_0}
 \longrightarrow\underline h_{(R_0,R_1,\ldots)}.
\]
Un \emph{cogrupoide en $\underline C$} es una cocategoría
$(R_0,R_1,\ldots)$ para la cual $h_{(R_0,R_1,\ldots)}(S)$ es un
grupoide para todo $S\in\underline C$; equivalentemente,
$\underline h_{(R_0,R_1,\ldots)}$ es una categoría cofibrada en
grupoides sobre $\underline C$.

En la práctica, las cocategorías y los cogrupoides aparecen del modo
siguiente. Sea $P:\underline F\longrightarrow\underline C$ una
categoría cofibrada, y fijemos un objeto $\xi$ de $\underline F$.
Pongamos $R_0=P(\xi)$, y denotemos por $(R_0,\underline C)$ la categoría
de flechas de $\underline C$ con fuente $R_0$. Sea
$P:\underline F'\longrightarrow(R_0,\underline C)$ la imagen inversa de
$\underline F$ por el funtor canónico
$(R_0,\underline C)\longrightarrow\underline C$. Elijamos una sección
cocartesiana $\xi^{\#}$ de $\underline F'$ sobre
$(R_0,\underline C)$ tal que
$\xi^{\#}(\operatorname{id}_{R_0})=\xi$. Así, para toda flecha
$R_0\xrightarrow fS$ en $\underline C$, elegimos un $f$-morfismo
cocartesiano $\xi\longrightarrow\xi^{\#}(f)$.

Definamos
\[
 \operatorname{Hom}_{\xi^{\#}}:
 (R_0\amalg R_0,\underline C)\longrightarrow(\mathrm{Sets})
\]
mediante
\[
 X=(f_1,f_2)\longmapsto
 \operatorname{Hom}_{\underline F(S)}
   (\xi^{\#}(f_1),\xi^{\#}(f_2)),
\]
donde $S$ es el blanco común de $f_1$ y $f_2$. Salvo equivalencia
natural canónica, este funtor no depende de la elección de $\xi^{\#}$;
en lo sucesivo escribiremos simplemente $\xi$ en lugar de $\xi^{\#}$.

Supongamos que el funtor está representado por un objeto
\[
 X_0=
 \left(R_0\sisi{\overset{s_1}{\underset{s_2}{=\!=}}}{\mathrel{\substack{\xrightarrow{\ s_1\ }\\[-.6ex]
                              \xrightarrow[\ s_2\ ]{}}}}R_1\right)
\]
de $(R_0\amalg R_0,\underline C)$. Existe entonces un $R_0$-morfismo
$\varphi:\xi(s_1)\longrightarrow\xi(s_2)$ en $\underline F$ tal que
\[
 \widetilde\varphi:
 \operatorname{Hom}_{(R_0\amalg R_0,\underline C)}(X_0,X)
 \longrightarrow\operatorname{Hom}_\xi(X)
\]
es biyectivo para todo $X$ en $(R_0\amalg R_0,\underline C)$. Entonces
$X_0$ define naturalmente una cocategoría $(R_0,R_1,\ldots)$ en
$\underline C$. En efecto, una terna de flechas

%%%%%%%%%%  scan idx 73  |  folio 62  %%%%%%%%%%

\[
\begin{tikzcd}[column sep=huge]
 R_0 \arrow[r,shift left=3,"u"]
     \arrow[r,"v"]
     \arrow[r,shift right=3,"w"']
 & S
\end{tikzcd}
\]
en $\underline C$ define una composición
\[
 \operatorname{Hom}_{\underline F(S)}(\xi(v),\xi(w))
 \times
 \operatorname{Hom}_{\underline F(S)}(\xi(u),\xi(v))
 \longrightarrow
 \operatorname{Hom}_{\underline F(S)}(\xi(u),\xi(w)),
\]
y, por tanto, una ley de composición
\[
 h_{R_1}(S)\times_{h_{R_0}(S)}h_{R_1}(S)
 \longrightarrow h_{R_1}(S),
\]
funtorial en $S$. Por consiguiente, $\operatorname{Hom}_\xi$ define la
cocategoría $(R_0,R_1,\ldots)$ en $\underline C$ y la categoría
cofibrada $\underline h_{(R_0,R_1,\ldots)}$ sobre $\underline C$.
Explícitamente,
\[
 \operatorname{ob}\underline h_{(R_0,R_1,\ldots)}
 =\coprod_{S\in\operatorname{ob}(\underline C)}
   \operatorname{Hom}_{\underline C}(R_0,S).
\]
Una flecha
\[
 (R_0\xrightarrow fS)\longrightarrow(R_0\xrightarrow{f'}S')
\]
es un par de flechas $S\xrightarrow hS'$ y $R_1\xrightarrow uS'$ en
$\underline C$ tales que
\[
 u\circ s_1=h\circ f,
 \qquad u\circ s_2=f'.
\]
El funtor estructural
$P:\underline h_{(R_0,R_1,\ldots)}\longrightarrow\underline C$ envía
$(R_0\xrightarrow fS)$ a $S$.

El par $(\xi,\varphi)$ define ahora un funtor
\[
 (\xi,\varphi):\underline h_{(R_0,R_1,\ldots)}
 \longrightarrow\underline F
 \quad\text{sobre }\underline C
\]
mediante
\[
 (R_0\xrightarrow fS)\longmapsto\xi(f)
\]
en los objetos y mediante
\[
 \left[(R_0\xrightarrow fS)
 \xrightarrow{(h,u)}(R_0\xrightarrow{f'}S')\right]
 \longmapsto
 \left[\xi(f)\xrightarrow{h_*}\xi(hf)
 \xrightarrow{u(\varphi)}\xi(f')\right]
\]
en las flechas. Para todo $S\in\underline C$, el funtor
\[
 (\xi,\varphi)(S):
 \underline h_{(R_0,R_1,\ldots)}(S)\longrightarrow\underline F(S)
\]
es plenamente fiel por la representabilidad de
$\operatorname{Hom}_\xi$. Se sigue que
\[
 (\xi,\varphi):\underline h_{(R_0,R_1,\ldots)}
 \longrightarrow\underline F
\]
es una $\underline C$-equivalencia si y solo si $\xi$ es cuasi-inicial
en $\underline F$ sobre $\underline C$, es decir, si todo objeto $a$ de
$\underline F$ admite una flecha cocartesiana
$\xi\longrightarrow a$.

También se puede considerar el funtor
\[
 \operatorname{Isom}_\xi:
 (R_0\amalg R_0,\underline C)\longrightarrow(\mathrm{Sets}),
\]

%%%%%%%%%%  scan idx 74  |  folio 63  %%%%%%%%%%

dado por
\[
 X=(f_1,f_2)\longmapsto
 \operatorname{Isom}_{\underline F(S)}(\xi(f_1),\xi(f_2)),
\]
donde $S$ es el blanco común de $f_1$ y $f_2$. Si este funtor está
representado por un objeto $(R_0\rightrightarrows R_1)$ de
$(R_0\amalg R_0,\underline C)$, entonces define un cogrupoide
$(R_0,R_1,\ldots)$ en $\underline C$, junto con un funtor
\[
 (\xi,\varphi):\underline h_{(R_0,R_1,\ldots)}
 \longrightarrow\underline F.
\]
Cuando $\underline F$ es un grupoide cofibrado, este funtor es una
$\underline C$-equivalencia si y solo si $\xi$ es cuasi-inicial en
$\underline F$ sobre $\underline C$.
